\section{Conclusion}\label{sec:conclusion}

We have shown that \EC{} between public and private consumption, combined with
empirically motivated heterogeneity in how constrained and unconstrained
households internalize publicly provided goods, transforms three features of
fiscal transmission. First, a new sufficient statistic -- the direct fiscal
channel $\xi$ -- emerges alongside the standard income-cyclicality statistic
$\chi$, giving researchers two independently-identifiable parameters to
discipline fiscal multipliers from data. Second, the curvature heterogeneity
induced by heterogeneous internalization decouples the aggregate MPC from the
aggregate MPE, providing a new route to the MPC/MPE trilemma resolution that
operates through the government-spending margin and is orthogonal to existing
labor-supply-based resolutions. Third, a distributional dominance condition
links inequality, complementarity, and multiplier size: fiscal policy is most
effective precisely in the economies where it is most needed -- those with
high inequality and strong complementarity between public services and the
private consumption of the poor.

Quantitatively, the HANK implementation with empirically disciplined
$\theta^{j}$ delivers impact multipliers in the range $1.68$--$1.84$, compared
to $1.50$ in the separable benchmark -- an economically significant amplification
that accounts for a meaningful share of the dispersion in the empirical fiscal
multiplier literature. Policy experiments reveal that targeted public spending
on complementary services can raise multipliers by another 24\% relative to
untargeted spending, while austerity implemented through cuts in public
services has aggregate contractionary effects 44\% larger than the separable
benchmark predicts -- with sharply asymmetric distributional costs falling on
the bottom of the distribution. In a developing-economy calibration with
higher $\lambda$ and stronger $|\theta^{H}|$, the impact multiplier rises to
$2.64$, consistent with the empirical evidence that fiscal policy is more
effective in low-income economies.

These findings have four broader implications. First, the fiscal-multiplier
literature's tradition of treating fiscal expansions as a single aggregate
object misses a first-order margin: the \emph{composition} of fiscal spending
across categories with different degrees of complementarity with private
consumption. Second, the welfare costs of austerity are systematically
under-estimated by models that suppress the $c$--$g$ complementarity channel,
and these costs fall disproportionately on constrained households. Third,
automatic stabilizers that maintain public-service provision during downturns
do ``double duty'': they stabilize aggregate demand \emph{and} protect the
private-consumption complementarity of the most vulnerable. Fourth, the
framework provides a micro-founded explanation for the reduced-form finding
that fiscal multipliers are larger in developing economies, and motivates
scaling up fiscal expansion through public services as a poverty-reduction
tool.

Several directions for future work emerge. \emph{Empirically}, the
household-level identification strategy of Section~\ref{sec:empirical} can be
extended to additional public-good categories and to the relationship between
$\theta^{j}$ and non-homothetic preferences emphasized by
\citet{BarthelFrancois2025}. \emph{Theoretically}, combining the $c$--$g$
non-separability of this paper with the $c$--$n$ non-separability of \BHL{}
would yield a framework with maximum flexibility to match all three MPC/MPE
facts simultaneously; we leave the cross-product analysis to future work.
\emph{Quantitatively}, the model can be extended to study transmission of
monetary policy under heterogeneous $\theta^{j}$, the interaction between
fiscal transmission and monetary normalization at the effective lower bound,
and the distributional consequences of fiscal consolidation in European and
emerging-market settings.

Most fundamentally, our results suggest that the integration of public goods
into household preferences -- long recognized as important but rarely
operationalized in macroeconomic models -- is a productive avenue for
macroeconomic research. When households care about publicly provided goods
differently depending on their income, the distributional incidence of public
spending becomes a first-order macroeconomic variable rather than a
distributional side-effect. This is a reframing that aligns the macroeconomic
analysis of fiscal policy with the growing evidence from development
economics, public finance, and inequality research that public-good provision
has shaped the distribution of economic outcomes at the global level.
